% part: intuitionistic-logic
% chapter: propositions-as-types
% section: proofs-to-terms

\documentclass[../../../include/open-logic-section]{subfiles}

\begin{document}

\olfileid{int}{pty}{pt}

\olsection{Ngowahi \usetoken{P}{derivation} dadi Terma Bukti}

Proses ngowahi !!{proof}s dedhuksi alami dadi terma bukti, yaiku
nerjemahake !!{proof}s ing~\Log{N2i} dadi !!{proof}s ing~\Log{tN2i},
prasaja banget. Kita mung nulis terma bukti ing sisih kiwa saben
!!{formula} ing sisih tengen sekuen ing !!{derivation}, supaya
ngasilake sekuen awujud $\Gamma \Sequent M : !A$. Kita banjur
ngarani $M$ \emph{nyekseni}~$!A$ ing konteks~$\Gamma$. Terma bukti
sing kudu ditulis ditemtokake sakabehe dening aturan \Log{tN2i}.

\begin{prop}
Saben !!{proof}~$\delta$ saka $\Gamma \Sequent !A$ ing \Log{N2i}
cocog karo !!a{proof}~$\delta'$ saka $\Gamma \Sequent N : !A$ ing
\Log{tN2i}. Terma bukti $N$ ditemtokake kanthi tunggal dening~$\delta$.
\end{prop}

\begin{proof}
Kita nggunakake induksi marang dhuwure~$\delta$. Yen $\delta$ iku
aksioma $x:!A \Sequent !A$, $\delta'$ iku aksioma~$x:!A \Sequent
x:!A$. Yen $\delta$ dipungkasi kanthi aturan inferensi, hipotesis
induksi menehi terma bukti kanggo saben premis, lan !!{proof}s
\Log{tN2i} saka sekuen sing ngandhut terma bukti kasebut lan cocog
karo saben premis. Kita ngetrapake aturan \Log{tN2i} sing cocog;
terma buktine ditemtokake dening aturan inferensi \Log{tN2i} kasebut.
\end{proof}

\begin{ex}
  Gatekna !!{proof} prasaja banget ing \Log{N1i} saka $(!A \land !B)
  \lif (!B \land !A)$:
  \[
  \AxiomC{$\Discharge{!A\land !B}{x}$}
  \RightLabel{\Elim\land[2]}
  \UnaryInfC{$!B$}
  \AxiomC{$\Discharge{!A\land !B}{x}$}
  \RightLabel{\Elim\land[1]}
  \UnaryInfC{$!A$}
  \RightLabel{\Intro\land}
  \BinaryInfC{$!B \land !A$}
  \DischargeRule{\Intro\lif}{x}
  \UnaryInfC{$(!A \land !B) \lif (!B \land !A)$}
  \DisplayProof
  \]
  Ing \Log{N2i}, wujude kaya mangkene:
  \[
  \Axiom$x : !A\land !B \fCenter !A\land !B$
  \RightLabel{\Elim\land[2]}
  \UnaryInf$x : !A\land !B \fCenter !B$
  \Axiom$x : !A\land !B \fCenter !A\land !B$
  \RightLabel{\Elim\land[1]}
  \UnaryInf$x : !A\land !B \fCenter !A$
  \RightLabel{\Intro\land}
  \BinaryInf$x : !A\land !B \fCenter !B \land !A$
  \RightLabel{\Intro\lif}
  \UnaryInf$\fCenter (!A \land !B) \lif (!B \land !A)$
  \DisplayProof
  \]
  Kita masang terma bukti saka ndhuwur menyang ngisor. Terma bukti
  kanggo aksioma padha karo variabel ing konteks, yaiku~$x$.
  Terma bukti sing digandhengake karo $\Elim\land[i]$ banjur
  ditemtokake minangka $\proj{2}{x}$ lan~$\proj{1}{x}$.
  Nalika ngetrapake \Intro{\land}, kita nggabungake rong terma
  bukti kasebut: $\pair{\proj{2}{x}}{\proj{1}{x}}$.
  Pungkasan, aturan \Intro{\lif} ngetrapake abstraksi $\lambd$,
  ngiket~$x$ lan nyathet yen $x$ menehi label asumsi~$!A \land
  !B$: $\lambd[\typeof{x}{!A \land
  !B}][\pair{\proj{2}{x}}{\proj{1}{x}}]$.
  Mula !!{proof} ing \Log{tN2i} iku:
  \[
  \Axiom$x : !A\land !B \fCenter x: !A\land !B$
  \RightLabel{\Elim\land[2]}
  \UnaryInf$x : !A\land !B \fCenter \proj{2}{x} : !B$
  \Axiom$x : !A\land !B \fCenter x: !A\land !B$
  \RightLabel{\Elim\land[1]}
  \UnaryInf$x : !A\land !B \fCenter \proj{1}{x} : !A$
  \RightLabel{\Intro\land}
  \BinaryInf$x : !A\land !B \fCenter \pair{\proj{2}{x}}{\proj{1}{x}} : 
    !B \land !A$
  \RightLabel{\Intro\lif}
  \UnaryInf$\fCenter \lambd[\typeof{x}{!A \land
  !B}][\pair{\proj{2}{x}}{\proj{1}{x}}] : (!A \land !B) \lif (!B \land !A)$
  \DisplayProof
  \]
\end{ex}


\end{document}
